% Lösungen Zone (zusammenbau v0.1)
\einheitenkopf[][Kennst du schon]{Das kennst du schon}

Zuordnung: 1 → Grundlage der Bernoulli-Kette und der Formel, Einheit 1 und 2, das Gegenereignis auch in Einheit 3 und 4 · 2 → für die Bernoulli-Bedingung in Einheit 1 · 3 → Anzahl der Anordnungen in der Bernoulli-Formel, Einheit 2, und beim Lesen von Summentermen in Einheit 3 · 4 → Werkzeug für Einheit 2 bis 5, Prüfungsteil B · 5 → Einheit 3 und 4 · 6 → Einheit 3 und 5 · 7 → Einheit 4 und die Aussagen über die Abhängigkeit von n · 8 → Lage des Maximums und Mitte einer Umgebung, Einheit 4 und 5 · 9 → Werkzeug für Einheit 2 bis 5, Prüfungsteil B · 10 → Werkzeug für Einheit 2 bis 5, Prüfungsteil B

\erg{1}{a) $0{,}2 \cdot 0{,}2 = 0{,}04$ \quad b) $0{,}6^2 + 0{,}4^2 = 0{,}36 + 0{,}16 = 0{,}52$}
\erg{2}{a) $\frac{4}{10} = 0{,}4$ – wie beim ersten Zug \quad b) $\frac{5}{20} \cdot \frac{4}{19} = \frac{1}{19}$}
\erg{3}{a) $4! = 24$ \quad b) $120$ Auswahlen}
\erg{4}{a) $0{,}65$ \quad b) $1 - 0{,}95 = 0{,}05$}
\erg{5}{a) $15$ Gäste \quad b) $11$ Personen ($0{,}15 \cdot 70 = 10{,}5$, aufrunden)}
\erg{6}{a) $34$ Familien \quad b) $16 + 6 = 22$ Familien}
\erg{7}{a) $0{,}125$ \quad b) $n = \frac{\mathrm{ln}\, 0{,}1}{\mathrm{ln}\, 0{,}6} \approx 4{,}51$}
\erg{8}{a) $60 \cdot \frac{1}{6} = 10$-mal \quad b) $240 \cdot 0{,}15 = 36$ Kunden}
\erg{9}{a) Fehler: „mindestens 3“ ist das Gegenteil von „höchstens 2“, nicht von „höchstens 3“. Richtig: $1 - 0{,}42 = 0{,}58$}
\erg{10}{a) $1 - 0{,}83 = 0{,}17$ (Gegenteil von höchstens 4)}
